\documentclass[11pt]{article}
\usepackage[utf8]{inputenc}
\usepackage{amsmath,amssymb}
\usepackage{graphicx}
\usepackage{booktabs}
\usepackage{authblk}
\usepackage[margin=1in]{geometry}
\usepackage[hidelinks]{hyperref}
\usepackage{setspace}

\newcommand{\grad}{\nabla}
\newcommand{\half}{\tfrac12}


\title{On Conduction by Hopping in RealQM:\\ \large the Free Boundary as Carrier}

\author[1]{Claes Johnson}
\affil[1]{KTH Royal Institute of Technology, SE-100 44 Stockholm, Sweden \\
\texttt{claesjohnson@gmail.com}}
\date{\today \\[0.4em]
\small\textit{Mathematical formulation by the author; analysis developed} \\
\small\textit{with AI assistance (Claude, Anthropic).}}

\begin{document}
\maketitle

\begin{abstract}
\noindent RealQM models matter as non-overlapping unit charge densities in ordinary
three-dimensional space, one per electron, meeting at free boundaries and minimising the Coulomb
energy. This article asks whether such a construction \textbf{conducts}.

What moves is not an electron. An electron here is a unit of charge density on a domain; what
exists is the \emph{partition} of space, and the carrier is a \textbf{free boundary} --- in the
relation a dislocation bears to the slip it produces, charge arriving at the far end while nothing
material crosses the sample. Single occupancy makes the construction the Hubbard model at
$U\to\infty$, so metals are out of scope in kind and not by degree, their carriers being extended
Bloch states; doped Mott insulators, molecular semiconductors and polaronic oxides are not.

The activation energy for that boundary to advance one site would decide the question, and it is
computable in principle: two configurations of a doped chain are fixed by \emph{symmetry} --- the
carrier centred on a site, and centred between two --- so both are stationary points at which the
interface condition holds without being imposed, and their difference is the hopping barrier, the
width of the slip being minimised over rather than chosen. What that measurement yields is
\textbf{a sign and a decade but not a value}. Transport is activated, the carrier sits at a
definite registry and must leave it to advance, and the cost is of order a tenth to a whole
electron-volt: every configuration computed, over four meshes, three chain lengths, four path
widths and two dimensions, falls between $68$ and $929$~meV. Within that range the grid decides.
The barrier is a difference of about $0.01$~Ha between two totals each carrying a discretisation
error twenty-five times larger, the cancellation such a difference needs does not occur, and
refining the mesh alone moves it $315\to552\to265$~meV without converging. Since the
$10$--$500$~meV window of real hopping semiconductors \emph{overlaps} that range, the construction
is neither placed in that class nor excluded from it. What it resembles is small-polaron or ionic
conduction rather than a metal. A resolution-limited null result is the honest outcome, and the
convergence study that establishes it is reported in full, since it bounds what any calculation of
this kind can currently claim.

A second obstacle is computational, and it conditions everything else. A free boundary
evolves by \emph{local} rules and so descends to the nearest configuration no single cell change
improves --- a stuck point, not a minimum --- and which one depends on where the run started. Four
independent geometries spread $0.005$--$0.05$~Ha across seeds and builds, against the
$0.001$--$0.02$~Ha barriers sought; the ambiguity exceeds the quantity by one to two orders, and is
unrepaired by surface tension, by step size, by freezing the labels, and by changing system. It is a
failure of the \emph{search}, not of the energy functional. The repair is to specify the partition
rather than find it: place the domain walls, freeze them, relax only the wavefunctions. Repeated
runs then agree to all five decimals recorded. On that footing a misplaced domain wall costs
$179$~meV, independent of chain length to $5\%$; a chain carrying a vacancy gains $0.112$~Ha by
delocalising its mismatch, also length-independent; and the barrier to sliding a pattern through its
lattice is tunable through zero by density and by core radius. The activation energy for a kink to
\emph{move}, which is what governs conductivity, is not among them and is left open with a stated
route rather than estimated.

Two further results bound the framework. The interface coefficient vanishes exactly, $C=0$, removing
degeneracy pressure and leaving a vanishing screening length and a dispersionless plasmon; a Robin
condition at $\beta a=1.146$ restores $C=C_{\rm TF}$, but $\beta$ that large \emph{unbinds}
$\mathrm{H}_2$, which fails by $\beta a\approx0.71$ --- so the equation of state and molecular
binding appear not to be repairable together, while double occupancy, the repair conduction needs
at half filling, becomes favourable at $\beta a=0.451$. And against density-functional
theory and exact full configuration interaction on the same chains, the framework errs by
\textbf{over-localising} --- in bond length and in pinning alike --- which is the direction
$U\to\infty$ implies, so it fails as its own analysis says it should rather than arbitrarily.
\end{abstract}

\setstretch{1.15}

\section{The extension, and the object it is made for}

In this article we extend RealQM, developed for atoms and molecules, to \textbf{two-component
systems}: ions and electrons, both charged, both free to move, interacting by Coulomb forces and
nothing else. A metal is such a system, and the question pursued here is the elementary one such a
system poses --- whether it \emph{conducts}, and if so what carries the charge. The same
construction covers a stellar interior and the recombination era of a Coulomb cosmology, which are
treated as the classical plasma, where $\omega_p^2=4\pi n$ comes out exactly.

RealQM \cite{realqm,realqmbook} represents an $N$-electron system by $N$ non-negative densities
$\psi_i$, each carrying unit charge on its own domain $\Omega_i$, the domains non-overlapping and
their interfaces free to move to where the energy is least. The energy minimised is Coulomb's:
\begin{equation}\label{eq:E}
E[\{\psi_i\},\{\Omega_i\}]=\sum_i\int_{\Omega_i}\Bigl(\half|\grad\psi_i|^2
-\sum_a\frac{Z_a}{|x-x_a|}\psi_i^2\Bigr)dx
+\half\sum_{i\neq j}\iint\frac{\psi_i^2(x)\psi_j^2(y)}{|x-y|}\,dx\,dy,
\end{equation}
subject to $\int_{\Omega_i}\psi_i^2=1$. There is no wavefunction in $3N$ dimensions, no
antisymmetry and no self-interaction; the role played by the exclusion principle is played by the
requirement that domains do not overlap. And --- of particular relevance here --- there is no smeared
continuum anywhere in the construction: everything in it is a discrete unit charge occupying a region
of ordinary space.

A two-component system of ions and electrons is therefore not an idealisation the theory must be
adapted to describe. It is what the theory is already made of, in a new arrangement.

Every previous test of this construction --- ionisation energies across the periodic table, molecular
binding, nuclear binding by dual Coulomb confinement \cite{realnucleus} --- has been on \emph{bound}
matter, with an attractor always present. A plasma is the first case in which large numbers of
charges of both signs are free to arrange themselves, and it is a fair question what the framework
then gives.

\section{What kind of conductor this can be}\label{sec:scope}

The construction does not have a free choice of target, and identifying it correctly settles most
of what follows. One electron per domain, non-overlapping, is usually described as the Hubbard model
at $U\to\infty$: double occupancy forbidden, every site singly occupied. That is right as far as it
goes and it is not far enough. The Hubbard Hamiltonian has two parameters, and non-overlap fixes
both. $U=\infty$ because two electrons may not share a domain; but also $t=0$, because the transfer
integral \emph{is} the overlap of neighbouring orbitals and there is none. The construction is
therefore not the Mott limit but the \textbf{atomic limit}, and the difference is not a refinement:

\begin{center}
\small
\begin{tabular}{@{}lll@{}}
\toprule
 & $U\to\infty$, $t\neq0$ & $U\to\infty$, $t=0$\\
\midrule
half filling      & Mott insulator      & isolated atoms\\
superexchange     & $J=4t^2/U$          & none\\
doped             & \textbf{conducts}   & still nothing\\
spin dynamics     & yes                 & none\\
\bottomrule
\end{tabular}
\end{center}

\noindent A Mott insulator is insulating but alive: it carries magnetic exchange, and doping it
produces a conductor, which is what the cuprates are. The atomic limit is a collection of atoms with
electrostatics between them. This accounts for every negative result below, and for one in
particular: doping a Mott insulator works \emph{through} $t$, so at $t=0$ there is nothing for
doping to activate --- which is exactly what the measurement shows, the barrier rising with carrier
concentration rather than falling. Metals are therefore out of scope in kind and
not by degree. Their carriers are extended Bloch states, their transport is band motion, and their
defining measurements --- de Haas--van Alphen oscillations, a linear electronic specific heat,
positive plasmon dispersion --- belong to a degenerate electron gas with a sharp Fermi surface,
which no arrangement of densities on domains appears able to supply. Two results usually read as
failures follow from the scope rather than against it: $d\sigma/dT>0$ is \emph{correct} for an
activated conductor and fatal only for a metal, and the absence of a Fermi surface does not bear on
a dilute carrier gas, which is non-degenerate and has none to reproduce.

What remains within reach is matter whose carrier is a localised charge on a site, whose gap is
on-site repulsion rather than band structure, and whose transport is activated hopping. That is a
real class --- doped Mott insulators, organic and molecular semiconductors, polaronic oxides,
amorphous semiconductors --- and this article set out to claim it. \S\ref{sec:conduction} measures
the activation energy and finds $613$--$1723$~meV, against the $10$--$500$~meV those materials
show. The claim does not survive its own measurement, and what the construction describes is the
insulating end of the class and not the conducting one. Silicon and the III--Vs were never in
scope, their gaps being band gaps.

\section{The interface, and the one number that is wrong}\label{sec:coefficient}

The free interface has a defect of its own, recorded here because it conditions everything the
partition does, though the hopping barrier does not turn on it. Partitioning a gas of density $n$
gives a kinetic energy density $t(n)=C\,n^{5/3}$ --- the Thomas--Fermi \emph{form}, obtained from
partition geometry rather than Fermi statistics --- but the flat trial function is admissible on a
free boundary, so $C=0$ exactly. There is no degeneracy pressure, the screening length vanishes
rather than being finite, and the plasmon does not disperse. The defect is one number and not the
construction: $C$ is fixed by the interface condition, and a Robin condition at $\beta a=1.146$
gives $C=2.871=C_{\rm TF}$ exactly, against Dirichlet's $14.804$, all three carrying the correct
$n^{5/3}$ scaling.

Whether that value is tolerable for matter is the open question, and it is sharper than a parameter
check. Adding the surface term and scanning it, $\mathrm{H}_2$ unbinds by $\beta a\approx0.71$,
below the $1.146$ the electron gas requires: Teller's theorem, that Thomas--Fermi binds no molecule,
is met inside the framework rather than inherited from outside it. Double occupancy --- the repair
conduction would need at half filling --- becomes favourable at $\beta a=0.451$ and is affordable by
comparison. These figures were obtained at differing densities and are not strictly commensurable;
the equation of state is pursued elsewhere, and what matters here is only that the interface
carries no cost of its own unless one is supplied.

\paragraph{A caution about specified partitions, learned by getting it wrong.} The method of
\S\ref{sec:conduction} evaluates the energy of a \emph{stated} partition, and it is sound only when
the stated partition is one the physics admits. The interface is not a surface to be shaped at
will: the wavefunction is a single global field partitioned by labels, continuous across every
interface, each domain carrying one unit of charge, and the boundary sits where the densities
balance. Freezing a partition removes that last condition. An attempt to measure a bending rigidity
by imposing a curvature, freezing it and relaxing the fields gave an energy falling symmetrically
and quadratically in the imposed curvature, which reads convincingly as a negative rigidity; but
released, a boundary started from the curved configuration returns to the flat one. The flat
interface is the balance surface and is stable, and what the frozen scan measured was the energy
gained by \emph{dropping the balance condition}. A specified partition gives a legitimate energy
difference between configurations the physics admits --- as symmetry guarantees for the two in
\S\ref{sec:conduction} --- and a meaningless one along a direction it forbids. Nothing in the
numbers distinguishes the two cases; the check is to release the boundary and see whether it stays.

\section{Computing with it anyway}\label{sec:conduction}

Electric conduction is the interaction of the electrons with the ion lattice. It is the elementary
case --- ordinary matter at ordinary temperature --- and it is for that reason the \emph{demanding}
one here, since ordinary-temperature conduction is degenerate matter, and degeneracy is what this
framework has least of.

One property of the construction is needed before anything else, and it is the property that makes
conduction a live question rather than a hopeless one. With the free interface used throughout, a
uniform density is admissible on any partition: $\psi_i\equiv|\Omega_i|^{-1/2}$ satisfies the
normalisation, minimises the electrostatic energy against a neutralising background, and has
$\grad\psi_i=0$. Redrawing the boundary therefore costs nothing,
\begin{equation}\label{eq:Czero}
\text{a free interface between domains of equal density carries no confinement energy,}
\end{equation}
and the boundary is the only object this framework moves --- its solver evolves each interface by the
local density balance $(\rho_i-\rho_j)/(\rho_i+\rho_j)$, transferring ownership of a cell when the
neighbour's density exceeds its own. Domains do not translate; they are redrawn. \emph{A vanishing
interface cost is a vanishing cost of transport.} What that same fact does to the equation of state
is a separate matter and is taken up in \S\ref{sec:coefficient}, where it is much less welcome.

Metallic conduction is the paradigm of \emph{delocalisation}: the carriers are extended states
spanning the crystal, whereas RealQM's electrons are localised by construction, each in its own
domain, with non-overlap requiring $N$ carriers to occupy $N$ disjoint regions. The natural minimiser
of \eqref{eq:E} is a set of compact cells, which describes a Mott insulator more naturally than a
conductor. Nor is there an evident mechanism for the metal--insulator distinction, which in standard
theory is band filling --- two electrons per band per cell --- an exclusion-principle count the
framework does not have; and with no Fermi surface, the carrier density in $\sigma=ne^2\tau/m$ has no
referent.

What does fit naturally is the opposite regime: localised charges hopping between domains by
thermally activated jumps, as in semiconductors, disordered systems and small-polaron transport. That
needs no Fermi surface and follows from the same thermal motion that supplies the pressure.

This yields a prediction crude enough to be decisive, and it concerns only a sign. Metallic
conductivity \emph{falls} with temperature; hopping conductivity \emph{rises}. If RealQM transports
charge by hopping rather than band motion it should give $d\sigma/dT>0$. Within the scope set in
\S\ref{sec:scope} that is the \emph{right} answer, activated transport being what a localised-carrier
semiconductor does; it would be fatal only for a metal, which is why the target matters. Whether it does has not been computed for electrons, and the
pessimistic reading should be held loosely: the same lattice that reproduces alkali lattice constants is
present here too.

For \emph{protons} the mechanism has already been demonstrated in this framework. Grotthuss
conduction --- an excess proton migrating along a hydrogen-bonded chain of water molecules by
successive transfers, rather than by bodily motion of any one particle --- is hopping transport in
its purest form, and it is a standard RealQM calculation: a chain of waters with $\mathrm{O}$--$\mathrm{O}$
spacing near $2.65$~\AA{}, an excess proton, and the question whether it hops from one molecule to
the next. Related transfers, $\mathrm{HF}+\mathrm{H_2O}\to\mathrm{F^-}+\mathrm{H_3O^+}$ and
$\mathrm{HCl}+\mathrm{NH_3}\to\mathrm{Cl^-}+\mathrm{NH_4^+}$, are computed the same way.

This matters for the argument rather than merely illustrating it. Charge transport by transfer
between localised domains is not a regime the framework must be stretched to reach --- it is one it
already handles, in the case where the carrier is a proton.

The electron case cannot presently be tested in the same way, and the obstruction is instructive
enough to be worth setting out. Two chain calculations were attempted here --- five hydrogens with an
excess electron, and five hydrogen kernels with one electron missing, the electronic counterparts of
the proton wire --- and both returned nothing. Neither result means what it appears to.

The cause is not what it first appeared. Kernel proximity assigns the labels only \emph{initially};
thereafter a free-boundary rule evolves each interface from the local density balance and transfers
ownership of a cell when the neighbouring domain's density exceeds its own. Partitions therefore do
rearrange, and electron transfer is representable in principle. What defeated the attempts here was
specific to each: a carrier given no kernel feels no force from the lattice; a bare kernel is given no
domain, so a vacancy has nothing to migrate into; and a carrier placed on an occupied site can seize
the nuclear region outright, since with free interfaces no confinement cost opposes it. Each is a
property of the construction, not of \eqref{eq:E}.

The deeper cause survives that repair. Conduction is charge in motion, and \eqref{eq:E} is an energy
minimisation: it returns the equilibrium partition, not a current. Given complete freedom in the
labels one would still obtain the ground state, and for a symmetric chain that state places one
domain per site by symmetry, with nothing flowing. A minimisation has no time in it.

\paragraph{Which filling the measurements belong to, and which needs the repair.} It is worth being
exact about this, because it determines how much of the article depends on $\beta$. At \emph{half
filling} --- one electron per site, every domain occupied --- the $U\to\infty$ construction forbids
conduction outright: there is no vacancy for a carrier to enter, and the insulator is structural
rather than dynamical. That is the case the double-occupancy repair is for. But it is not the case
measured here. The chains reported above are doped by ionising every second atom, leaving half the
sites empty, and every barrier in this article is the cost of a carrier moving into a \emph{vacancy}.
Single occupancy handles that without amendment.

So the transport results stand on their own: they require no finite $U$, no double occupancy and no
$\beta$. What requires the repair is the undoped, half-filled chain, which is not run here. This is
the stronger position of the two --- the measurements are independent of the proposal made to extend
them --- and it locates precisely what the proposal would buy: not better hopping, but conduction at
a filling where the construction currently permits none.

What the framework does possess is an asymmetry of dynamics. Nuclei move, under classical equations
of motion; electrons relax, to a static minimum. Charge carried by \emph{nuclei} is therefore
representable --- which is precisely why the Grotthuss calculation works, the carrier there being a
proton --- and charge carried by electrons is not, because in this scheme electrons have no equation
of motion at all. The correct machinery exists elsewhere in the programme, in the time-dependent form
where a real density carries a complex current, and a conduction calculation belongs there rather
than here.

\paragraph{What the barrier means as a field.} A barrier becomes a threshold field when
divided by the charge and the hop distance. Taking the range of \S\ref{sec:mesh} rather than any
single value, $E_a/ea$ runs from $0.07$ to $0.9$~V per lattice site at $a=6.0$; copper carrying an
ampere runs at $4\times10^{-12}$~V per site. The comparison spans an order of magnitude and is
uncertain by that much, but ten orders separate the two either way, so the conclusion --- that no
field a material survives will drive this by direct depinning --- does not depend on where in the
range the barrier falls.

\paragraph{Why a chain rather than a ring.} A ring supplies a clean barrier --- no ends,
every site equivalent, the energy obliged to repeat with the lattice period --- but it cannot host
the drive: a conservative field has zero circulation, so no bias drives a sea around a closed loop.
A chain has ends, a potential drop along it is an ordinary thing to apply, and a single carrier may
enter at one end and leave at the other, which is the transport event itself.

\paragraph{The partition must be specified, not searched for.} Before the number, the method,
because the obvious approach fails and the failure is instructive. A kink is \emph{not} the ground
state, so a free-boundary minimisation relaxes it away --- asked to hold one, the solver oscillates
between the kinked and uniform partitions and the energy never settles, wandering over
$0.05$~Ha against a barrier of order $0.01$. That is not a convergence failure to be tuned away:
it is unaltered by the surface tension $\beta$, by reducing the boundary step fivefold, and by
freezing the labels outright.

The cause is general and it disqualifies the free boundary for every transport quantity in this
article. The interface moves by \emph{local single-cell label flips}, so a run descends to whichever
configuration no single flip improves --- the nearest stuck point, not the minimum --- and which one
that is depends on the initial condition. Four independent geometries agree:

\begin{center}
\small
\begin{tabular}{@{}lll@{}}
\toprule
geometry & spread & against a barrier of\\
\midrule
ring, between two builds & $0.005$~Ha & $0.012$--$0.022$\\
kink chain, across steps & $0.05$~Ha (never settles) & $\sim0.01$\\
$\mathrm{H}^-+\mathrm{H}$, four seeds & $0.031$~Ha & $\sim0.001$\\
alkali chain, two seeds & $0.013$~Ha & $\sim0.001$\\
\bottomrule
\end{tabular}
\end{center}

\noindent In every case the ambiguity exceeds the quantity sought, by one to two orders. It is a
failure of the \emph{search}, not of the energy functional, which is why nothing in the functional
touches it. It is also why the ring's $0.01186$~Ha moved to $0.022$~Ha under a recompilation that
altered no physics: rounding at the last bit flips cells sitting near the threshold, and a different
stuck point follows. The rigid-sliding figure should be read as a bound of that order rather than a
determination, and the doping barriers of the earlier barrier scans, obtained the same way, with it.

The repair is not a better minimiser but no minimiser. The domain walls of a chain are planes
perpendicular to the axis, so they can be \emph{placed}: specify the partition, freeze it, and relax
only the wavefunctions. Nothing is searched for, the same input returns the same output, and the
question asked is no longer ``where does the boundary go'' but ``what does this configuration
cost'' --- which is the question a barrier actually poses. It is the device that makes the
$\mathrm{H}_2$ midplane exact, applied to a coordinate instead of a symmetry.

\paragraph{The hopping barrier, from two configurations neither of which was chosen.} The
quantity conduction depends on is the activation energy for a carrier to advance one site, and it
can be had without choosing any partition, because \emph{symmetry fixes both ends}. Take a chain
with an odd number of sites carrying one vacancy, and let the carrier hop between the two central
sites. Two configurations are then distinguished:

\begin{itemize}\itemsep2pt
\item the carrier centred \textbf{on} a site --- invariant under reflection about that site;
\item the carrier centred \textbf{between} the two sites --- invariant under reflection about the
bond midpoint.
\end{itemize}

\noindent Each is invariant under a reflection, hence a stationary point of the energy over
partitions, hence a configuration at which the balance condition holds without being imposed. This
is the footing that makes the $\mathrm{H}_2$ midplane exact, and it is the only footing on which a
specified partition is legitimate. The profile connecting the two is given the same width in both,
so that position is the only thing that differs.

\begin{center}
\small
\begin{tabular}{@{}lrrr@{}}
\toprule
spacing $a$ (a$_0$) & on site (Ha) & between sites (Ha) & barrier\\
\midrule
$4.0$ & $-1.00761$ & $-1.07092$ & $0.06331$~Ha $=1723$~meV\\
$5.0$ & $-1.14835$ & $-1.17695$ & $0.02860$~Ha $=\phantom{0}778$~meV\\
$6.0$ & $-1.21951$ & $-1.24957$ & $0.03006$~Ha $=\phantom{0}818$~meV\\
$7.0$ & $-1.17275$ & $-1.19527$ & $0.02252$~Ha $=\phantom{0}613$~meV\\
\bottomrule
\end{tabular}
\end{center}

\paragraph{And the barrier must be minimised over the path, not read off one.} A barrier is the
lowest saddle connecting two states, not the highest point of whichever path one happens to sample.
The family above has a second parameter --- the width $W$ over which the slip is spread --- and
fixing it describes a kink of rigid shape, which a minimum-energy path is under no obligation to be.
Scanning it:

\begin{center}
\small
\begin{tabular}{@{}lrl@{}}
\toprule
$W$ (sites) & barrier & \\
\midrule
$0.4$ & $929$~meV & \\
$0.8$ & $818$~meV & \\
$1.5$ & $552$~meV & minimum over $W$\\
$3.0$ & $626$~meV & \\
\bottomrule
\end{tabular}
\end{center}

\noindent The curve has a minimum, which is what a Frenkel--Kontorova kink requires: elastic
stiffness penalises a wide kink, the lattice's grip penalises a narrow one, and the physical width
lies between. Every figure obtained at a fixed $W$ is an upper bound, since it is the cost along a
rigid path; the $818$~meV that this section reported before the width was varied is one such.

These four figures share one mesh, $h=0.348$, and the minimum over $W$ is therefore a minimum at
that mesh. \S\ref{sec:mesh} varies $h$ and finds that the barrier does not survive it. The width
scan is reported here because the shape of the curve --- a genuine minimum, for the reason
Frenkel--Kontorova gives --- is what one wants to know about the path family, and that shape is not
what fails. The number attached to it is.

\noindent Three things follow, and the first is a structural claim rather than a number, though a
narrower one than it first appeared. On a chain the carrier sits \textbf{between} sites: the
bond-centred configuration is the minimum at every spacing tested, so the hop is bond to bond
through the on-site saddle, which reproduces by an independent route the registry preference found
above. That preference does \emph{not} survive a second dimension. Repeating the same construction
on a square lattice --- identical carrier, identical slip, one dimension more --- the hole sits on a
\emph{site} and the bond becomes the saddle, which is the conventional picture for vacancy hopping.
The bond-centred result is a property of the chain and not of the construction, and it is recorded
here as such.

Second, the barrier looked like a property of the hop rather than of the chain: at $a=6.0$ and
$W=0.8$ it is $0.03006$~Ha for seven sites and $0.03247$~Ha for nine, $8\%$ apart, while the total
energy moves by $0.59$~Ha between them. \S\ref{sec:mesh} withdraws this. The two figures were
computed at different meshes, and at matched mesh the same comparison spreads by a factor of five.

Third, \textbf{compression raises the barrier}. It is
$818$~meV at the equilibrium spacing and $1723$~meV compressed to $a=4.0$, a doubling rather than a
closing. That is the conclusion the compression scans of this article reached by the free-boundary
route, and it is worth stating plainly that those scans are the ones \S\ref{sec:conduction} shows to
be unreliable. The result survives the discrediting of the method that first found it.

\paragraph{Nor does doping, which was the last candidate.} Carrier concentration was the one
variable with evidence of moving the barrier downward: the free-boundary series of this article has
$1.09$~eV at $29\%$ falling to $0.31$~eV at $43\%$, a factor of $3.5$. Repeating the measurement
with symmetry-fixed endpoints --- two vacancies in a seven-site chain, one hopping while the other
stays --- gives $1525$~meV at $29\%$ against $818$~meV at $14\%$. The barrier \emph{rises}, and the
earlier trend is not reproduced. It belongs with the other figures obtained by that route.

One qualification belongs with the comparison. The carrier's preferred position changes with
concentration: at $14\%$ it sits at a bond and the site is the saddle, at $29\%$ the reverse. The
two hops are therefore not the same process, and $818\to1525$ should not be read as a clean trend.
What matters for the conclusion is that both lie between one and one and a half electron-volts.

\paragraph{Two dimensions raise it further.} The chain forces the carrier along one route, and
one might expect a lattice to offer a cheaper one --- an earlier scan of this framework found an
\emph{excess electron} threading between two atoms to cost $5.46$~eV against $31.5$~eV through one,
a factor of $5.8$. A hole has no such option. Built identically on a square lattice, the barrier is
$2072$~meV at $3\times3$ and $2911$~meV at $5\times5$; the smaller lattice has a single interior
site and is dominated by its edges, which is also where the bond-centred preference of the chain
survives and then fails. Taking the larger figure, two dimensions make the hop \emph{harder} by a
factor of $3.6$, not easier: a missing electron at a site disturbs four neighbours rather than two,
and there is no route that avoids them.

\paragraph{The conclusion is therefore robust rather than provisional.} Four variables have been
tried, each of which could have produced a hopping semiconductor:

\begin{center}
\small
\begin{tabular}{@{}ll@{}}
\toprule
variable & effect on the barrier\\
\midrule
lattice spacing, compressed & doubles it, $818\to1723$~meV\\
a second dimension          & triples it, $818\to2911$~meV\\
carrier concentration       & doubles it, $818\to1525$~meV\\
core radius                 & helps modestly\\
\bottomrule
\end{tabular}
\end{center}

\noindent The barrier remains between $0.56$ and $2.9$~eV throughout, that is $21$ to
$113\,k_BT$, against the $10$--$500$~meV of the materials this article hoped to describe. It is not
one measurement that places the construction outside that class but four attempts to bring it
inside, none of which succeeds, and two of which make matters worse in the direction one would
naively expect to help. \emph{No accessible variation reduces the barrier to the range where
hopping occurs.}

\section{The mesh, and why no barrier is quoted}\label{sec:mesh}

Every barrier above is a difference of two total energies, each computed on a finite grid. Whether
such a difference means anything depends on whether the two discretisation errors cancel, and that
is a question about the mesh, not about the physics. It was asked late, and the answer withdraws the
numbers.

Holding chain length, spacing, core radius and path width fixed at $n=7$, $a=6.0$, $r_c=1.6$,
$W=1.5$, and varying only the grid:

\begin{center}
\small
\begin{tabular}{@{}lrrrr@{}}
\toprule
$h$ (a$_0$) & on site (Ha) & between sites (Ha) & barrier & slope of total\\
\midrule
$0.451$ & $-1.34727$ & $-1.35882$ & $315$~meV & \\
$0.397$ & $-1.31216$ & $-1.32578$ & $371$~meV & $0.650$\\
$0.348$ & $-1.26709$ & $-1.28739$ & $552$~meV & $0.920$\\
$0.299$ & $-1.22260$ & $-1.23234$ & $265$~meV & $0.908$\\
\bottomrule
\end{tabular}
\end{center}

\noindent \textbf{The barrier does not converge.} It rises, then falls below where it started:
$315\to371\to552\to265$~meV, a factor of $2.1$ across four meshes with no monotone trend and nothing
to extrapolate. The scatter is not noise --- three independent runs of the same configuration agree
to all five decimals recorded, as everywhere else in this article once the partition is specified.
It is exact, reproducible dependence on the grid.

The totals, by contrast, behave. Their increments settle to a constant slope in $h$ of about
$0.91$~Ha per unit $h$, which is first-order convergence, expected for a sharp-interface method.
That is what makes the failure legible: at $h=0.30$ the discretisation error in each total is still
about $0.27$~Ha, some twenty per cent of the total, while the barrier extracted from the pair is
$0.010$~Ha. \emph{The quantity sought is under four per cent of the error in each quantity it is
differenced from.} For it to survive, the two errors would have to cancel to better than
ninety-six per cent, and they do not: the on-site slope settles at $0.91$ while the between-sites
slope is still climbing through $0.61$, $0.78$, $1.12$. Extrapolating the two separately to $h=0$
gives a difference near $2$~eV, inconsistent with every raw value in the table --- which is another
way of exhibiting the same non-cancellation.

Chain length behaves no better once the mesh is matched, which retracts a claim made above. At
$h\approx0.45$, seven sites give $315$~meV and eleven give $560$; at $h\approx0.40$, seven give
$371$ and nine give $68$. The agreement between $552$ at $n=7$ and $560$ at $n=11$ that earlier
sections read as length-independence compared $h=0.348$ against $h=0.449$: two errors cancelling by
coincidence. \emph{No length-independence has been established.}

The same caveat reaches the other two length-independence claims of this article --- the
$179$~meV wall defect and the delocalisation gain of $0.112$~Ha, each reported constant to a few
per cent across chain length. Both were obtained the same way, at fixed requested spacing on a box
that grows with the chain, so the mesh varied with the length there too. Neither has been rechecked
at matched resolution. They are larger quantities than the hop barrier and so less exposed --- the
delocalisation gain by a factor of eleven --- but the constancy should be read as untested rather
than as confirmed.

\paragraph{What is therefore withdrawn, and what stands.} The figure $E_a=555$~meV is withdrawn.
So is $265$~meV, and every other entry in these tables taken as a measurement: they are samples from
a range $68$--$929$~meV set by the grid, and the calculation cannot distinguish among them. The
earlier paragraphs of this section, which compare barriers across spacing, dimension and carrier
concentration, all rest on differences of the same size computed at a single mesh, and inherit the
same limitation; their \emph{trends} may well be right, since the configurations compared are alike
and their errors more likely to cancel, but no weight can be put on the factors of two.

Two things do stand. The \textbf{sign} is robust: in every configuration computed --- four meshes,
three chain lengths, four widths, both dimensions --- the on-site configuration lies above the
bond-centred one on a chain. A barrier exists, and the carrier has a preferred registry it must
leave to advance. And the \textbf{decade} is robust: every value obtained lies between $68$ and
$929$~meV. That is not the meV scale of a metal, and not the ${\sim}13$~eV of the on-site repulsion
$U$; it is the tenth-of-an-electron-volt scale of small-polaron and ionic conduction. The framework
produces activated transport of that order.

What it cannot do is decide whether the barrier lies inside the $10$--$500$~meV window of hopping
semiconductors. That was to have been the verdict of this article, in one direction or the other,
and it is not available at this discretisation. Resolving it needs error cancellation the present
sharp-interface treatment does not provide --- a finer grid by a factor of two or more, a
higher-order interface, or a formulation in which the barrier is computed directly rather than as a
difference of large totals.

\paragraph{What the barrier says about scope.} Less than this article set out to establish. The
activation energy is of order a tenth to a whole electron-volt, every value obtained falling between
$68$ and $929$~meV, and the grid rather than the physics decides where in that range any particular
figure lands. Hopping semiconductors --- the organic and molecular solids, the polaronic oxides, the
amorphous semiconductors named in \S\ref{sec:scope} --- have activation energies of $10$ to
$500$~meV. That window overlaps the range computed here and is not separated from it, so the
comparison does not decide the question either way.

The scope claim of \S\ref{sec:scope} is therefore neither established nor refuted. What can be said
is that the construction gives \emph{activated} transport at the tenth-of-an-eV scale, with a
carrier that has a definite preferred registry and must leave it to advance --- the phenomenology of
small-polaron and ionic conduction rather than of a metal. Whether it also reaches the
semiconducting window is a question this discretisation is too coarse to answer, and it is left open
rather than decided by a number the mesh study does not support.

\paragraph{Earlier figures, and why they are not quoted.} A series of barrier scans was made
before the difficulty above was understood, all of them by the free-boundary route: a compression
test on a ring, a doping series at two concentrations, a lattice barrier between interstitial
pockets of order $5$~eV, and a localised interface defect costing $179$~meV independently of chain
length. Every one specified or evolved a partition that no symmetry fixes, so each is a bound of
its stated order rather than a determination, and all are superseded by the barrier above. What
survives from them is the direction of the density trend, which the symmetry-fixed measurement
confirms.

\paragraph{What is actually transported.} The carrier in this framework is not an electron. An
electron here is not a particle but a unit of charge density on a domain, and what exists is the
\emph{partition} of space; dynamics is the evolution of that partition, and transport is the
partition pattern travelling. What moves along the chain is the \textbf{free boundary} --- a locus
where two densities balance --- and nothing material traverses it at all.

The quantitative form of this is the useful one. As a kink passes, each electron advances by exactly
one site spacing and then stops, while the kink advances the entire length of the chain. The
carrier's displacement and its constituents' displacement are decoupled, and their ratio is the
chain length. This is the relation between a dislocation and the slip it produces: the dislocation
crosses the crystal, no atom moves far. Charge is genuinely transported and real charge arrives at
the far end; it is the \emph{carrier} that is immaterial, in the sense that a wave is immaterial
while the water is not.

That diagnosis also accounts for the failures recorded earlier in this study rather than excusing
them. Each of the unsuccessful constructions imposed a carrier that had to \emph{translate}, moving
bodily through occupied space, and returned costs of tens of electron-volts or else collapsed. Those
were the price of forbidding the only mechanism the framework has. The solver states as much
directly: its free boundary evolves by the density balance $(\rho_i - \rho_j)/(\rho_i + \rho_j)$,
flipping ownership cell by cell. Boundary motion is the only motion it performs. The corona result
is the same statement seen from the favourable side --- an added electron spread into a uniform
shell at no cost, which is a boundary redistributing freely.

This raises the stake of the kink measurement. If transport works here it works \emph{because} of
the free boundary, which is the single feature separating this framework from the standard theory;
and if the boundary cannot carry charge cheaply, the difficulty is not in a detail of the
construction but in the framework's defining object.

\paragraph{A repair that is not available, and one that is.} An obvious response to the barrier is
to let valence electrons share a territory, which removes it by construction. The $n^{5/3}$
scaling rules that out: sharing destroys the extensivity the partition supplies, and a metal would
lose its kinetic energy density by fifteen orders of magnitude.

\paragraph{What the framework then is.}\label{sec:hubbard} Non-overlapping sites, occupation $0$, $1$ or $2$, an on-site
repulsion $U$, and transfer between neighbours: that is the \textbf{Hubbard model}. Which locates
RealQM exactly. As formulated --- one unit of charge per domain, always --- it is the Hubbard model in
the limit $U\to\infty$, where double occupancy is forbidden outright and every site is singly
occupied. That limit is a Mott insulator at every density, which is precisely what
the earlier barrier scans found by direct computation, and it explains the density trend as well: no
finite $U$ can be beaten by a bandwidth that does not exist.

The repair is correspondingly minimal. Keep the partition, keep the geometry, keep the Coulomb
energetics, and allow occupation to take the value two at a computed cost. The framework then sits
where the standard treatment of strongly correlated matter sits --- the Hubbard model is the working
description of Mott insulators, of the cuprates, and of the materials for which band theory fails ---
with the difference that its $U$ and its lattice are derived rather than supplied.

We do not carry this out here. What is established is where the obstruction lies: not in the
interface condition, not in the geometry, but in the rule of one unit of charge per domain. That rule
is also the source of one of the framework's results, charge quantization as a consequence of
indivisibility rather than an assumption. The tension is worth stating plainly: \emph{the postulate
that explains why charge comes in units is the postulate that forbids a metal.} Matter has it both
ways, quantizing the particles while leaving the amplitude on each site continuous, and a partition
of real densities has no evident means to do the same.

\section{The same chain, computed the ordinary way}\label{sec:stdqm}

Everything above is internal to the framework. It is worth putting the same system beside the
standard treatment, on quantities both can produce, and the result is a check on the scope claim of
\S\ref{sec:scope} rather than a contest.

Kohn--Sham DFT, PBE/6-31G, all electrons, a straight chain of lithium atoms at the spacings used
above, with the geometry held uniform so that the comparison is like for like --- a Peierls
distortion is not permitted on either side:

\begin{center}
\small
\begin{tabular}{@{}lrrr@{}}
\toprule
 & RealQM & PBE/6-31G & experiment\\
\midrule
equilibrium spacing (a$_0$)      & $4.5$              & $6.0$    & $5.7$ (bcc n.n.)\\
verdict at $a=6.0$               & pinned, $184$~meV  & metallic & lithium conducts\\
verdict at its own equilibrium   & depinned, $\approx0$ & metallic & ---\\
cohesive energy (eV/atom)        & $\sim1.7$          & $0.486$  & $1.63$ (bulk, 8 n.n.)\\
\bottomrule
\end{tabular}
\end{center}

\paragraph{And it is not one functional's answer.} Density-functional theory is not the
Schr\"odinger equation: it replaces the interacting problem with fictitious non-interacting
electrons in an effective potential, exact in principle by Hohenberg--Kohn but resting in practice
on a \emph{constructed} exchange--correlation functional and a finite basis --- two uncontrolled
approximations. The comparison above is therefore model against model, not model against exact, and
would normally be qualified accordingly. On the quantity at issue it need not be. Repeating the
scan with Hartree--Fock, which carries no correlation whatever, and with MP2, which carries it
perturbatively, moves the total energy by some $0.3$~Ha between methods and leaves the equilibrium
spacing at $6.0$~a$_0$ in all three:

\begin{center}
\small
\begin{tabular}{@{}lrrr@{}}
\toprule
$a$ (a$_0$) & HF & PBE & MP2\\
\midrule
$4.5$ & $-51.97126$ & $-52.29669$ & $-52.02861$\\
$5.0$ & $-52.01844$ & $-52.33386$ & $-52.06750$\\
$6.0$ & $\mathbf{-52.04726}$ & $\mathbf{-52.34332}$ & $\mathbf{-52.08273}$\\
$7.0$ & $-52.04089$ & $-52.31653$ & $-52.06476$\\
\bottomrule
\end{tabular}
\end{center}

\noindent The twenty per cent is thus a disagreement with standard quantum chemistry rather than
with a choice of functional.

\paragraph{But the objection cuts the other way, and harder.} If one may choose the approximation to
suit the problem then one will always win, and that is not a comparison. Applied evenly, the charge
falls on \emph{this} side of the table. The standard column used a non-empirical functional, a
general-purpose basis, and three treatments of correlation that were not selected after the fact ---
but the RealQM column carries a free parameter, the pseudopotential radius $r_c$, and it was set to
$1.6$~a$_0$ by choice rather than by any independent criterion. At $r_c=2.4$ the equilibrium spacing
of the same chain moves out beyond $6.0$~a$_0$, which \emph{agrees} with the standard result. The
twenty per cent is therefore a statement about $r_c=1.6$ and not about the framework, and the
framework could have been made to agree as easily as to disagree.

A comparison is only worth making if the parameters are fixed by something other than the quantity
being compared. The proper protocol is to calibrate $r_c$ on the \emph{isolated} lithium atom ---
requiring the single pseudo-atom to reproduce the measured first ionisation energy,
$0.198$~Ha --- and then to compute the chain with that value and no further adjustment. The chain is
then a prediction rather than a fit, and stands on the same footing as the standard column, in which
nothing was tuned to lithium chains. Until that is done the table above should be read as
establishing the \emph{method-independence of the standard answer} and nothing about the size of any
discrepancy.

\noindent The metallic verdict is not read off a single gap, which any finite chain has, but off its
scaling: $0.432$, $0.346$, $0.291$, $0.247$, $0.173$~eV at $n=5,7,9,11,13$, with $n\times{\rm gap}$
constant near $2.4$~eV, so the gap closes as $1/n$ and is finite-size rather than a band gap.

\paragraph{That verdict is weaker than it looks, and in a direction that matters.} What was compared
are Kohn--Sham eigenvalue differences, and Kohn--Sham eigenvalues are not excitation energies. They
belong to a fictitious system of \emph{non-interacting} particles introduced to reproduce a density;
the true gap differs from theirs by the derivative discontinuity of the exchange--correlation
functional, which PBE does not possess at all. Approximate functionals are known to underestimate
gaps for exactly this reason, often by a factor of two. A small Kohn--Sham gap is therefore what the
method returns whether or not a gap is present.

The bias has a direction, and it is the one at issue here. Approximate functionals carry a
self-interaction error --- an electron repelled in part by its own density --- whose consequence is
\textbf{delocalisation error}: densities spread too far, localised solutions are penalised, and
metallic answers are favoured. So the standard column calling this chain metallic is that treatment
erring in its characteristic direction, not a neutral measurement. Against it, the framework of this
article over-localises, by construction, being $U\to\infty$. \emph{Both are biased and in opposite
senses}, and lithium is simply a case where the standard bias does no damage, because the chain
really is metallic. On a system where it does damage --- an oxide, a doped Mott insulator, anything
whose gap is on-site repulsion --- the delocalisation error is the well-known failure and the
localisation is the correct instinct.

Two further asymmetries belong in the record. Density-functional theory admits no convergence path:
wavefunction methods form a hierarchy, Hartree--Fock through MP2, CCSD, CCSD(T) to full
configuration interaction, which converges to the exact answer within a basis, whereas no procedure
refines PBE towards the true functional. And the quantities are not of the same kind. What is
computed here are total energies of \emph{specified configurations} --- differences between states
of one system, which are observables --- where the gap above is an eigenvalue difference in an
auxiliary system that was never claimed to have physical states. Where this framework is less
accurate it is at least computing something that means what it says.

\paragraph{What this establishes, and what it does not.} On this system the standard treatment is
better and cheaply so: the equilibrium spacing within a few per cent of the measured
nearest-neighbour distance, the right conduction verdict, a sensible cohesive energy, in seconds.
RealQM places the atoms some twenty per cent too close, and at the spacing lithium actually adopts
it predicts a pinning barrier of $184$~meV where the metal has none.

That is not a surprise and should not be read as one, because \emph{lithium is outside the scope
this article claims}. \S\ref{sec:scope} excludes metals in kind and not by degree: extended Bloch
states, band transport, a Fermi surface, none of which the construction possesses. A simple metal is
the standard treatment's best case and this framework's declared worst, so the table is a scope check
and not a competition.

As a scope check it passes, in the specific sense that matters: the framework errs by
\emph{over-localising}, in bond length and in pinning alike, which is what $U\to\infty$ implies and
therefore the direction the identification predicts. It fails as it was said it would, rather than
arbitrarily.

\paragraph{Where the comparison would actually bite.} The corollary is uncomfortable and is recorded
plainly: nothing computed here is a quantity the standard treatment does not already supply more
accurately. The case for the construction rests instead on the real-space mechanism --- a registry,
a sliding partition, a carrier that is a boundary rather than a particle --- being a more transparent
account of activated hopping than a model Hamiltonian with fitted parameters. That claim is not
tested by lithium. It would be tested on the systems the scope table names, where on-site repulsion
rather than band structure sets the gap and where density-functional theory is itself in difficulty:
a doped Mott insulator, an organic semiconductor, a polaronic oxide. Those comparisons are not made
here.

\section*{A note on what was tested numerically}

The plasma frequency $\omega_p^2=4\pi n$, the screening relation, the Thomas--Fermi comparison and
the vanishing of $C$ are analytic. The alkali lattice constants are a radial Wigner--Seitz
calculation with the stated pseudopotential and correlation functional, run for five alkalis; the
standard column reproducing the
series is what makes the comparison meaningful, and its own $44\%$ mean error in the bulk modulus
sets the accuracy available at this level of model. Three attempts are reported as null rather than omitted. The interface parameter
$\beta$ was tested on helium at three meshes at fixed imaginary time, giving shifts of $+0.151$,
$-0.018$ and $-0.184$~Ha as the mesh refined --- swinging through zero and growing, with absolute
energies moving away from the exact $-2.9037$. The cause is the geometry rather than the arithmetic:
in helium the two domains meet at a plane \emph{through the nucleus}, so interface error and
Coulomb-cusp error are superimposed and cannot be separated on a uniform grid. The measurement wants
$\mathrm{H_2}$, where the domains meet at the midplane \emph{between} the nuclei in a smooth
region. An interface-condition test on helium was
not converged, giving a sensitivity that reversed sign under mesh refinement, and no conclusion is
drawn from it. Two chain calculations --- five hydrogens with an excess electron, and five hydrogen
kernels with one electron missing --- returned no transport, but for a reason established by reading
the constructions rather than by the runs: a carrier without a kernel feels no lattice force, and a
bare kernel receives no domain for a vacancy to occupy. Both are recorded because the null results
located those obstructions. A caution on the numerics belongs with them. Several attempts returned
energies below the physical floor for their particle content, but these track the boundary speed and
the seeding rather than the configuration: with the interface driven faster than the densities can
follow, or with a carrier seeded directly onto an occupied kernel, the relaxation runs away, while the
same configurations at a quarter of that speed settled to physical values. A two-domain shell on one
nucleus is in fact stable under the density-balance rule --- compressing the inner domain raises its
density at the interface and lowers the outer's, which pushes the boundary back out --- and both
domains carry real kinetic energy there, since the nucleus shapes their profiles. The vanishing of
$C$ is a statement about a \emph{uniform} density and does not transfer to a shaped one. Scripts and pages accompany the source.



\section{Conclusion}\label{sec:conclusion}

The question was whether a construction of non-overlapping unit charge densities, meeting at free
boundaries, conducts. The answer is not a number, and the most useful thing this article can report
is why.

Transport here is activated: the carrier has a preferred registry and must leave it to advance,
and the cost of doing so is of order a tenth to a whole electron-volt --- every configuration
computed, across four meshes, three chain lengths, four path widths and two dimensions, falling
between $68$ and $929$~meV, and every one of them putting the same configuration on top. The sign
and the decade are solid. The value is not. It is a difference of order $0.01$~Ha taken between two
total energies each carrying a discretisation error some twenty-five times larger, and
\S\ref{sec:mesh} shows the required cancellation does not occur: refining the grid alone moves the
barrier from $315$ to $552$ to $265$~meV without converging. An earlier draft of this article quoted
$E_a=555$~meV and drew a verdict from its position relative to $500$; that figure was one sample
from a grid-dependent range, and it is withdrawn.

The consequence is that the question the article set out to settle stays open. The
$10$--$500$~meV window of hopping semiconductors overlaps the range computed here rather than
sitting outside it, so the construction is not shown to belong to that class and is not excluded
from it either. What it does resemble, on the evidence that survives, is small-polaron or ionic
conduction: activated, tenth-of-an-eV, with a carrier dragging a lattice distortion rather than
propagating freely.

Three things stand independently of the number, and they are of different kinds.

\paragraph{The carrier is a free boundary, and that picture holds.} An electron here is a unit of
charge density on a domain; what exists is the partition of space. Transport is therefore the
partition pattern travelling, in the relation a dislocation bears to slip: charge arrives at the far
end while nothing material crosses the sample. That diagnosis accounts for the null results reported
above, each of which had imposed a carrier obliged to translate bodily through occupied space, and
it was confirmed directly --- given a geometry that permits it, the boundary migrates and the charge
goes with it. The picture is not in question.

\paragraph{Costing it defeats a minimisation, and that is the article's principal result.} A free
boundary evolves by local rules, so a run descends to the nearest configuration no single cell
change improves. That is a stuck point and not a minimum, and which one is reached depends on the
initial condition. Measured across four independent geometries the ambiguity is $0.005$--$0.05$~Ha
against barriers of $0.001$--$0.02$: one to two orders larger than the quantity being sought,
unrepaired by the surface tension, by the boundary step, by freezing the labels, or by changing to a
system without the first one's pathologies. It is a failure of the search and not of the energy
functional, which is why nothing in the functional touches it, and it invalidates transport energies
obtained this way --- including, by its own logic, the rigid-sliding figure of this article, which
moved under a recompilation that altered no physics.

The repair is not a better minimiser but none. Domain walls are surfaces and can be placed:
specify the partition, freeze it, relax only the wavefunctions, and ask what a stated configuration
costs rather than where a boundary goes. Repeated runs then agree to every decimal recorded. What
that yields is modest but real --- a misplaced domain wall at $179$~meV, independent of chain length
to $5\%$; a vacancy's mismatch gaining $0.112$~Ha by delocalising, also length-independent; a
sliding barrier tunable through zero by density and by core radius. What it does not yet yield is
the activation energy for a kink to \emph{move}, which is the quantity conductivity depends on, and
which is left open here with a construction stated rather than a number estimated.

\paragraph{What the postulate buys, and what it costs.} The results here say something about
non-overlap itself, and it is worth stating separately from any of them. Overlap is not an optional
feature of quantum mechanics; it is the mechanism by which amplitude passes between places.
Forbidding it makes the electrons objects with disjoint supports --- distinguishable, in the way
classical charges are --- and leaves a theory of \emph{static} charge distributions carrying a
quantum-derived kinetic penalty.

That division predicts where the framework succeeds and where it does not, and the prediction is
borne out across its whole corpus. Ground-state energies, bond lengths, lattice constants, binding
curves: quantities set by electrostatics against a localisation cost, and it does them well.
Tunnelling, exchange, dispersion, transport, a Fermi surface: quantities that exist only because
amplitude is shared, and it has none of them. The failure recorded in this article is not a
shortfall in accuracy but an absence of the ingredient, and it could have been predicted from the
postulate without any computation --- which is the strongest form of the result, and the least
comfortable.

\paragraph{The construction is $U\to\infty$, and this sets both its scope and its errors.} Single
occupancy is the Hubbard limit in which double occupancy is forbidden, so metals are excluded in
kind --- no extended states, no band transport, no Fermi surface --- and what remains in scope is
matter whose gap is on-site repulsion. The same fact fixes the direction of every discrepancy
measured against standard treatments: the framework \emph{over-localises}, in bond length and in
pinning alike. It fails as its own analysis says it should, which is worth more than an accidental
agreement, but it does not yet produce a quantity that established methods do not supply more
accurately. Its case rests on the mechanism being a clearer account of activated hopping than a
model Hamiltonian with fitted parameters --- a claim tested not by simple metals, where
density-functional theory is at its best, but by the systems where on-site repulsion sets the gap
and that theory is itself in difficulty.

Two repairs were identified and priced. Finite $U$ in place of infinite, which conduction at half
filling requires, becomes favourable at $\beta a=0.451$ and leaves molecules bound. Restoring the
interface coefficient to its Thomas--Fermi value requires $\beta a=1.146$, and $\mathrm{H}_2$
unbinds by $\beta a\approx0.71$. Teller's theorem is thereby met inside the framework rather than
inherited from outside it.

\begin{thebibliography}{9}
\bibitem{realqm} C.~Johnson, \emph{Real Quantum Mechanics}, \texttt{https://realqm.com}.
\bibitem{realqmbook} C.~Johnson, \emph{Real Quantum Mechanics: a new view of the atom}, book
manuscript.
\bibitem{realnucleus} C.~Johnson, \emph{RealNucleus: nucleons as charge densities}, submitted.
\bibitem{madelung} E.~Madelung, \emph{Quantentheorie in hydrodynamischer Form},
Z.~Phys.~\textbf{40} (1927) 322.
\end{thebibliography}

\end{document}
